% note éditée à la main (coupe du 17/11/2026) : graphiques.fin ne l'écrase pas
DGFiP, REI 2017 et 2025, DVF 2023 à 2025 ; Insee, recensement 2022, comptes de patrimoine 2025 et Insee Analyses n° 91 ; calculs de l'auteur (synthese/calculs\_foncier.py, foncier/simulation.py).\par Note : chaque variante remplace la moitié des droits sur les logements par de la taxe foncière, à recette constante et sans réaction des ventes ; la proposition remplace aussi la moitié des droits sur les autres immeubles. Taux départemental : chaque département remplace ses propres droits par un taux sur le sol des logements de son territoire.
